\chapter{Misconceptions Taxonomy}
In order to systematically analyze the misconceptions surrounding \gls{llms}, and the effects of the developed mediation tool, we have developed a taxonomy of misconceptions. This taxonomy is based on a review of the literature on \gls{llm} misconceptions, as well as an analysis of my own experiences and observations of \gls{llm} misconceptions. The field of \gls{llm} misconceptions is vast and rapidly evolving, but still its early stages, and there is no widely accepted taxonomy of misconceptions.

Taxonomies are classification systemts aimed at helping to organize and understand complex phenomena based on characteristics.\cite{kundisch_update_2022} Their use is widely spread in numerous fields to provide a framework for classification, idenfication and description of phenomena.\cite{schoormann_exploring_2022}

Since the field of \gls{llm} misconceptions is still in its early stages, we have developed a taxonomy that is based on our own analysis and observations, rather than relying on existing taxonomies. Due to the vastness of the field, we have divided the taxonomy into four dimensions, each of which captures a different aspect of \gls{llm} misconceptions. These dimensions are: foundational conceptual accuracy, mechanistic understanding, ethical and academic integrity, and socio-cognitive impact. Each dimension is further divided into subcategories, which capture specific types of misconceptions within that dimension.


\section{Foundational Conceptual Accuracy}
\begin{table}[htbp]
\centering
\caption{Dimension 1: Foundational Conceptual Accuracy}
\label{dim1}
\begin{tabularx}{\textwidth}{|l|l|X|}
\hline
\textbf{Misconception} & \textbf{Label} & \textbf{Description} \\
\hline
Anthropomorphism & ANTHRO & Attributing human mental states, intentions, or understanding to \gls{llms}. \\
\hline
Deterministic Automation & DETAUTO & Confusing \gls{llms} with deterministic automation systems. \\
\hline
Accuracy Overestimation & OVERACC & Overestimating the reliability and correctness of \gls{llm} outputs. \\
\hline
Autonomy Overestimation & OVERAUTO & Belief that the \gls{llm} operates independently or possesses agency. \\
\hline
\end{tabularx}
\end{table}

This dimension of the taxonomy is focused on misconceptions related to the foundational conceptual understanding of \gls{llms}, and also the confusion between \gls{llms} and other types of AI systems. The term ``AI'' exploded in popularity after the relase of GPT-3.5 model, and the term is often used to refer to a wide range of technologies, including \gls{llms}. This has lead to a lot of confusion and misconceptions about what \gls{llms} are and what they can do.

\subsection{Anthropomorphism}
The category of anthropomorphism is defined as the attribution of human mental states, intentions, or understanding to \gls{llms}. This misconception seems to be common amongst users who interact with \gls{llms} in a conversational manner, like chatbots. Subjects was placed in this category by using terms such as ``understand'', ``think'', ``learn'' and ``know'' to describe the \gls{llm}'s capabilities in a non-metaphorical way.

Anthropomorphism emerges both from the system design of \gls{llms}, and from the user's interpretation of the system's behavior.\cite{maeda_when_2024} A paper by Maeda and Quan-Haase provides a framework for when \gls{hai} become parasocial. They argue that chatbot's variety in information presentation resemble day-to-day bi-directional communication, and that the users' tendency to fill in missing conversational context causes them to attribute human-like qualities to the system. The roleplaying nature of the chatbot interface is visible through language choices such as personal pronouns and paraphrasing suggests attributes such as intention and understanding in the system which lead users to project inference to the system due to the lack of information about the system's inner workings.\cite{maeda_when_2024} 
This is empirically supported by a large scale study by Ibrahim et al.\ which found similar anthropomorphic behaviours in all tested \gls{llms} (Gemini 1.5 Pro, Claude 3.5 Sonnet, GPT-4o and Mistral Large).\cite{ibrahim_multi-turn_2026} They found, in particular, use of personal pronouns, vallidation and empathy as the most common anthropomorphic behaviours, and that these behaviours mostly occured after multiple turns of conversation.

The anthropomorphism misconception is widespread. An investigation by Colombatto and Fleming in 2024 found that in a sample of 300 US residents, only a third of the participants thought that ChatGPT definetly did not have subjective experiences, while two thirds applied varying degrees of phenomenal consciousness.\cite{colombatto_folk_2024} The investigation also showed that the participants' conciusness attributions increased with usage frequency, rather than correcting the misconception.

Building on this, Lin\cite{lin_six_2025} conceptualizes anthropomorphism as an ``anthropomorphsm fallacy'', in which users and researchers misinterpret \gls{llm} outputs as evidence of underlying mental life. The fluent, human-seeming responses of \gls{llms} can lead to an enhanced ELIZA effect, encouraging users to attribute understanding, beliefs, or intentions to systems that in reality operate through statistical token prediction. This creates a ``language-user illusion'', where the system's linguistic capabilities are mistaken for genuine cognition. 

Furthermore, the anthropomorphic metaphores used surroundning \gls{llms} can reinforce this misconception.\cite{lin_six_2025} Unlike humans who are influenced by their own experiences, \gls{llms} simulate a wide range of viewpoints based on their training data\cite{lin_six_2025}, which can lead to overestimation of their capabilities.

\subsection{Deterministic Automation}

\subsection{Accuracy Overestimation}

\subsection{Autonomy Overestimation}

\section{Mechanistic Understanding}

\section{Ethical and Academic Integrity}

\section{Socio-Cognitive Impact}
